% Source: ME34_v1.html
\documentclass[11pt]{article}
\usepackage[utf8]{inputenc}
\usepackage[T1]{fontenc}
\usepackage[margin=1in]{geometry}
\usepackage{amsmath}
\usepackage{amssymb}
\usepackage{siunitx}
\usepackage{enumitem}
\usepackage{parskip}

\title{PS160 -- ME34\_v1.html}
\date{}

\begin{document}
\maketitle

\section*{Questions}

\begin{enumerate}[leftmargin=*]
  \item An object is placed 21 cm to the left of a flat mirror.

\textbf{(a)} What is the position of the image? (Positive on the same side as the object; negative on the opposite side. Answer with an integer.) [BLANK]
cm

\textbf{(b)} What is the magnification of the image? (Answer with an integer.) [BLANK]

\textbf{(c)} State whether the image is real or virtual. [BLANK]

\textbf{(d)} State whether the image is upright or inverted. [BLANK] \hfill \textit{(10 pts, Fill-in-blank)}
    \begin{itemize}
      \item -21
      \item 1
      \item virtual
      \item upright
    \end{itemize}

  \item A spherical mirror is concave and has focal length 10.7 cm. If the object distance is 23 cm, calculate the image distance, \emph{s}', in cm. \hfill \textit{(10 pts, Numeric)}

  \item An eggplant is 1.3 cm from the vertex of a convex spherical mirror having a radius of curvature of 5.8 cm. Calculate the lateral magnification. \hfill \textit{(10 pts, Numeric)}

  \item An object is located 10 cm in front of a converging lens of focal length 14.0 cm. The image point is [DROPDOWN: -19 | -35 | -28 | -22]
cm, the magnification is [DROPDOWN: 2.5 | -2.5 | 3.5 | -3.5 | 1.9]
, and the image is both [DROPDOWN: upright | inverted]
and [DROPDOWN: real | virtual]
. \hfill \textit{(10 pts, Dropdown)}

  \item A thin diverging lens has a focal length of 9.00 cm. If the object distance is 16.0 cm, the image distance is [DROPDOWN: 8.36 | -8.36 | -12 | -5.76 | 12 | 5.76]
cm, the magnification is [DROPDOWN: 0.42 | 0.36 | -0.42 | 0.68 | -0.36 | -0.68]
, and the image is both [DROPDOWN: virtual | real]
and [DROPDOWN: upright | inverted]
. \hfill \textit{(10 pts, Dropdown)}

  \item A lens is made of a material of index 2 with R1 = 8 cm and R2 = -11 cm. If the object is at 23 cm, what is the image position (in cm)? \hfill \textit{(10 pts, Numeric)}

  \item A moon pebble embedded in a transparent, spherical material with index of refraction 2 has an image that is inside the sphere and appears to be 14 cm from the surface. How far from the surface is the actual object? The radius of curvature of the sphere is 33 cm.

(Hint: What kind of image is it? Pay very close attention to the sign conventions!) \hfill \textit{(10 pts, Numeric)}

  \item The figure (not to scale) shows a tiny tree 21 cm from a converging lens followed by a diverging lens with \emph{L}= 10 cm. Find the \emph{x-}position in cm of the final image if \emph{f}\textsubscript{1} = 8 cm and |\emph{f}\textsubscript{2}| = 17 cm.

[Image: two\_lenses\_02.png] \hfill \textit{(10 pts, Numeric)}

  \item A woman has a far point of 43.3 cm. What power lens in diopters does she require so that she can see distant objects? \hfill \textit{(10 pts, Numeric)}

  \item What focal length in centimeters is required for a magnifier if the magnification is to be 2.3? \hfill \textit{(10 pts, Numeric)}

  \item A [DROPDOWN: converging | diverging]
lens corrects for near-sightedness. The image it creates is [DROPDOWN: real | virtual]
. A [DROPDOWN: diverging | converging]
lens corrects for far-sightedness. The image it creates is [DROPDOWN: virtual | real]
. \hfill \textit{(10 pts, Dropdown)}

  \item For a spherical mirror, if the object is farther away than the focal point, then the image is [DROPDOWN: virtual | real]
and [DROPDOWN: upright | inverted]
. If the object is closer than the focal point, then the image is [DROPDOWN: real | virtual]
and [DROPDOWN: upright | inverted]
. \hfill \textit{(10 pts, Dropdown)}

\end{enumerate}

\end{document}